\chapter{Discussion and Lessons Learned}
\label{ch:discussion}

This chapter interprets the results so far and collects the lessons recorded in the project's learning log
(\code{docs/learning_log.md}). The project is also a study in harness engineering, so the lessons are part of the
contribution: most of them are about how small decisions in the harness change what a cheap model does.

\section{What the results suggest about the research question}

The research question asks whether a graph-specific harness lets a cheap model match a frontier model in a
general-purpose harness at a fraction of the cost. The evidence so far points in that direction on simple tasks and
is silent on hard ones.

On the ten dataset problems the cheap model, with fitting tools, was as accurate as the frontier model in Claude Code
(9 of 9 each in the preview, 270 of 270 lenient in the M2 run), while the general-purpose harness used 17--35 times
more input tokens per question. On these tasks the outcome is decided by networkx in both systems, so the comparison
is essentially a comparison of overhead, and the overhead of a coding harness is large. This is the cost side of the
research question in its simplest form.

The accuracy side is still open, for three reasons. The dataset is saturated: on large graphs every task is at 100\%
lenient accuracy, so it cannot show where a cheap model falls behind. The harder tasks so far (waypoint, combine) were
run on 3--15 questions, once, with one model at a time. And the fair comparison, in which both harnesses use the same
model and the same tools, has not been run. The combine results do show something the dataset could not: that a cheap
model in our harness can fail badly (2 of 15 with tools alone) and that a harness-level change (offering code with a
strategy hint) can recover most of it (14 of 15). Whether a frontier model in a general harness would do better on the
same tasks, and at what cost, is the M6 question.

\section{Lessons on tools}

\paragraph{A fitting tool beats a checker.} Three times a tool that answers the task directly (\code{has_edge},
\code{degree}, \code{distances_from} with \code{neighborhood}) kept accuracy and cut time or tokens by roughly an order
of magnitude. When such a tool exists, a small model stops reasoning about how to answer. A fitting tool removes an
error instead of catching it, which is cheaper than any verifier. The design consequence is to move decisions the
model is bad at (what ``degree'' means here, counting a 45-item list) into tools.

\paragraph{Building blocks, not one-off tools.} The multi-step tools were not written for the probe questions; they
answer a family of distance and neighborhood questions. The opposite extreme, one tool per question, would make the
tool list grow without bound, which conflicts with showing only relevant tools.

\paragraph{Tool descriptions are the model's manual, and the schema wins.} When the instructions said ``submit the
handle'' and the schema said ``array of edges'', the model followed the schema. When \code{has_path}'s description
promised only ``a path'', the model still used it as a shortest path, relying on a property the tool did not promise.
Anything the model may do must be allowed by the schema, and anything it relies on should be stated in the
description.

\paragraph{Overlapping tools create ambiguity.} \code{has_path} and \code{shortest_path} overlap; the model used both
for the same purpose. Either the overlap is documented or it is removed; per-family tool filtering (M4) reduces it.

\paragraph{An interface the model knows beats a better one it has to learn.} Inside code, tool functions that returned
our JSON dictionaries failed every question; plain values in the style of networkx worked.

\paragraph{Return evidence.} A tool should return the proof of its answer (the path, the cycle, the coloring or the
odd cycle), so a verifier can check the answer later without solving the question again. This is also the main
argument for keeping tools when code is more accurate: a number printed by code carries no evidence.

\section{Lessons on the loop and the prompt}

\paragraph{The harness can cause hallucinations.} A generic nudge (``use the tools or submit'') pushed the model from
an honest ``my tools are not enough'' into a guess. The answer schema shapes behaviour: without an honest exit, the
model submitted 0 for a node that did not exist. \code{cannot_answer} fixed both, and introduced its own risk: it can
become an easy way out on questions that the tools could answer. Its rate should be tracked as a metric of its own.

\paragraph{Format, not reasoning, is where the small model fails now.} On the dataset, the errors left are tool calls
written as text and evidence left out of a ``yes''. The harness turns many of these into correct answers (lenient
mode), which is part of what makes a small model usable, and which is why both lenient and strict accuracy are
reported. Result handles may reduce the problem on long-list tasks, since the model can submit a handle instead of
writing out 45 nodes; this has not been measured yet.

\paragraph{Facts the model needs for planning belong in the context.} The node-degree slowdown came from thinking that
happens before any tool call: the model debated whether the graph was directed. The fact was one tool call away, but
a tool cannot help with a question that is asked before the first call. Facts that do not give away an answer
(directedness, yes; node count, no) belong in the prompt up front. This is the open fix B.

\paragraph{Check every field before re-sending it.} The model's hidden thinking was being re-sent on every step,
more than tripling the prompt of the second step. What goes back into the history matters as much as what goes into the prompt.

\paragraph{Nothing the model reads may grow with the graph.} Result handles covered tool results, but error messages
and repeated values still scaled with the graph (a single typo produced 59,000 characters on a 10,000-node graph). A
happy-path test on a big graph does not exercise the error paths; every message needs the same rule.

\paragraph{One hint for every task is too blunt.} The code strategy hint was decisive for qwen3.5 on the combine tasks.
In the invalidated first round it also appeared to hurt qwen3:8b (the model guessed a dictionary key and failed on it
repeatedly), but that round was affected by the context limit and cannot be used as evidence. Either way, different
models need different guidance, which is the case for per-family skills (M4) and for a skill ablation (none, minimal,
prescriptive) per family and per model. The GDS Agent's history points the other way for frontier models: its skill
was cut from 109 lines to 22 within a week.

\section{Lessons on verification}

\paragraph{A verifier checks evidence; it never re-solves.} A verifier that recomputes the answer with networkx and
compares is an oracle, and makes any comparison meaningless. This restricts what can be checked
(Table~\ref{tab:checkable}): counts, degrees and every ``no'' are unverified. The shortest-path verifier's length check
is a borderline case (it computes the optimal distance) and should be replaced by a certificate check.

\paragraph{Code verifiers are better than a model judge, where they apply.} A verifier in code gives a precise,
trustworthy error message (``2--5 is not an edge in G'') and costs nothing; a second model as a judge would be
neither precise nor free. This is the same pattern coding agents use with tests and compilers.

\paragraph{Verifiers bought nothing on the dataset.} Because the tools already computed the right answer, there was
nothing to catch. Verifiers should matter where answers are assembled by the model from several results (waypoint
routes) or produced by code, and where escalation (M5) can act on a failed check. That remains to be shown.

\section{Lessons on method}

\paragraph{Check the serving settings before blaming the model.} A silent 4,096-token context produced behaviour that
looked exactly like a model weakness (going silent after finding the answer). The context window of every local run
is now logged and reported.

\paragraph{Read the failures, not only the score.} ``The model failed'' was, on inspection, ``the question was
ambiguous'' (the waypoint wording) or ``the schema forbade the right answer'' (handles).

\paragraph{One ideal sequence is the wrong yardstick for process evaluation.} A run that read the way back from an
earlier result instead of making a second call was correct and well reasoned. Process metrics need an ``alternative
path'' class, and judging whether each step is justified is more useful than matching one expected sequence.

\paragraph{Repeated runs matter.} The same question about a missing node gave different outcomes across four runs.
One run is not a result; the plan's three runs and bootstrap intervals are a requirement, not a refinement.

\paragraph{The log is for analysis, not for plumbing.} Reading results back from the trace file made the runner
quadratic; the code that writes the log already knows the values.

\paragraph{Pin model identifiers.} The alias \code{sonnet} resolved to different model versions on different days.

\section{Limitations}

Beyond the limitations of the setup listed in Chapter~\ref{ch:setup}, two limitations of the approach should be
named. First, the harness is only as good as its tools: a question type without a fitting tool falls back to many
low-level calls or to code, and the combine results show how badly the first can go. Second, many interesting answers
cannot be verified from evidence alone, so a large part of any benchmark will be reported as unverified, and the
benefit of verification and escalation is limited to the tasks where evidence exists.
